% =====================================================================
\section{Mock Exam -- Form B}\label{ch:exam-b}
% =====================================================================

\textbf{Instructions.} Answer all problems. Show complete solutions; box your
final answers. Time: 2 hours. Total: 100 points.

\subsection*{Part I. Vector functions (35 points)}

\begin{questions}
\item Describe the curve $\vr(t)=\ip{t,\;\cos2t,\;\sin2t}$. \pts{5}
\item Evaluate $\displaystyle\lim_{t\to0}\ip{\frac{e^t-1}{t},\;\frac{\sqrt{1+t}-1}{t},\;\frac{3}{1+t}}$. \pts{5}
\item Find a vector equation and parametric equations of the line segment from
      $P(1,-2,3)$ to $Q(4,0,-1)$, and use it to find the midpoint. \pts{5}
\item Find the unit tangent vector $\vT(0)$ for
      $\vr(t)=\ip{e^{2t},\;e^{-2t},\;te^{2t}}$. \pts{6}
\item Reparametrize $\vr(t)=\ip{2t,\;1-3t,\;5+4t}$ with respect to arc length
      measured from the point where $t=0$, in the direction of increasing $t$. \pts{6}
\item For $\vr(t)=\ip{2\sin t,\;2\cos t,\;\sqrt5\,t}$ find $\vT(t)$, $\vN(t)$,
      $\vB(t)$ and $\kappa(t)$. \pts{8}
\end{questions}

\subsection*{Part II. Directional derivatives and the gradient (35 points)}

\begin{questions}[resume]
\item Use the \emph{limit definition} to find $\Du f(1,1)$ for $f(x,y)=x^2-y$ and
      $\vu=\ip{\frac35,\frac45}$. \pts{8}
\item Find $\Du f(-6,4)$ for $f(x,y)=\sin(2x+3y)$, where $\vu$ is the unit vector
      at angle $\theta=-\frac\pi6$. \pts{7}
\item Let $f(x,y,z)=x^2+2y^2+3z^2$.
      \begin{enumerate}[label=(\alph*)]
        \item Find the maximum rate of change of $f$ at $(4,-1,1)$ and its direction.
        \item Find the tangent plane and normal line to the ellipsoid
              $x^2+2y^2+3z^2=21$ at $(4,-1,1)$.
      \end{enumerate}\pts{12}
\item At a point $P$, $D_{\vi}f(P)=-3$ and $D_{\vj}f(P)=4$. Find $\grad f(P)$, the
      direction of fastest \emph{decrease} at $P$, and the rate in the
      direction of $\vi+\vj$. \pts{8}
\end{questions}

\subsection*{Part III. Vector calculus (30 points)}

\begin{questions}[resume]
\item Match each field with its description.\\
      (i) $\dfrac{\ip{-y,x}}{\sqrt{x^2+y^2}}$\quad (ii) $\ip{x,y}$\quad (iii) $\ip{0,x}$\quad (iv) $\ip{1,1}$\\[2pt]
      P. All arrows are equal.\\
      Q. Arrows are vertical: up in the right half-plane, down in the left.\\
      R. Arrows of length 1, tangent to circles about the origin, counterclockwise.\\
      S. Arrows point away from the origin and grow with distance. \pts{6}
\item Show that $\vF=\ip{e^x\sin y+2x,\;e^x\cos y-1}$ is conservative and find a
      potential function. \pts{8}
\item Find the potential function $f$ of
      $\vF=\ip{y^2z^3,\;2xyz^3,\;3xy^2z^2}$ with $f(1,1,1)=4$. \pts{8}
\item Is $\vF=\ip{y,\;x,\;xyz}$ conservative? Justify. If you try the
      potential-function method anyway, where does it break down? \pts{8}
\end{questions}

\begin{center}\textit{--- End of Form B ---}\end{center}

% =====================================================================
\subsection{Solutions to Form B}

\solhead{1}{5}
$y=\cos2t$ and $z=\sin2t$ satisfy $y^2+z^2=1$, while $x=t$ increases steadily.
The curve lies on the circular cylinder $y^2+z^2=1$ (axis = the $x$-axis): it is
a \textbf{helix about the $x$-axis} of radius $1$, making one full turn every
time $x$ increases by $\pi$ (the period of $\cos2t$). At $t=0$ it passes through
$(0,1,0)$.

\solhead{2}{5}
Componentwise: $\lim\frac{e^t-1}{t}=1$ (standard limit);
$\frac{\sqrt{1+t}-1}{t}=\frac{(1+t)-1}{t(\sqrt{1+t}+1)}=\frac{1}{\sqrt{1+t}+1}\to\frac12$;
$\frac{3}{1+t}\to3$. Answer: $\boxed{\ip{1,\frac12,3}}$.

\solhead{3}{5}
$\vr(t)=(1-t)\ip{1,-2,3}+t\ip{4,0,-1}=\ip{1+3t,\;-2+2t,\;3-4t}$, $0\le t\le1$;
i.e. $x=1+3t$, $y=-2+2t$, $z=3-4t$. Midpoint ($t=\frac12$):
$\boxed{\left(\frac52,-1,1\right)}$.

\solhead{4}{6}
$\vr\,'(t)=\ip{2e^{2t},\;-2e^{-2t},\;e^{2t}+2te^{2t}}$ (product rule in the third
component). $\vr\,'(0)=\ip{2,-2,1}$, $\norm{\vr\,'(0)}=3$, so
$\boxed{\vT(0)=\ip{\frac23,-\frac23,\frac13}}$.

\solhead{5}{6}
$\vr\,'(t)=\ip{2,-3,4}$, $\norm{\vr\,'}=\sqrt{29}$. $s(t)=\int_0^t\sqrt{29}\,\dd u=\sqrt{29}\,t$, so
$t=s/\sqrt{29}$ and
\[
  \boxed{\vr(t(s))=\ip{\frac{2s}{\sqrt{29}},\;1-\frac{3s}{\sqrt{29}},\;5+\frac{4s}{\sqrt{29}}}},\quad s\ge0 .
\]
(Check: the new velocity $\ip{2,-3,4}/\sqrt{29}$ has length $1$.\checkok)

\solhead{6}{8}
\step{1}{Speed.} $\vr\,'=\ip{2\cos t,\;-2\sin t,\;\sqrt5}$, $\norm{\vr\,'}=\sqrt{4+5}=3$.
\step{2}{$\vT$.} $\vT=\frac13\ip{2\cos t,\;-2\sin t,\;\sqrt5}$.
\step{3}{$\vN$.} $\vT\,'=\frac13\ip{-2\sin t,\;-2\cos t,\;0}$, $\norm{\vT\,'}=\frac23$, so
$\vN=\ip{-\sin t,\;-\cos t,\;0}$.
\step{4}{$\vB=\vT\times\vN$.}
\[
  \vB=\frac13\begin{vmatrix}\vi&\vj&\vk\\2\cos t&-2\sin t&\sqrt5\\-\sin t&-\cos t&0\end{vmatrix}
  =\frac13\ip{\sqrt5\cos t,\;-\sqrt5\sin t,\;-2\cos^2t-2\sin^2t}
  =\frac13\ip{\sqrt5\cos t,\;-\sqrt5\sin t,\;-2}.
\]
\step{5}{$\kappa$.} $\kappa=\dfrac{\norm{\vT\,'}}{\norm{\vr\,'}}=\dfrac{2/3}{3}=\boxed{\dfrac29}$
(constant, as for every helix: $\frac{a}{a^2+b^2}=\frac{2}{4+5}$).

\solhead{7}{8}
\[
\frac{f\left(1+\frac35h,1+\frac45h\right)-f(1,1)}{h}
=\frac{\left(1+\frac65h+\frac9{25}h^2\right)-1-\frac45h-0}{h}
=\frac{\frac25h+\frac9{25}h^2}{h}=\frac25+\frac{9}{25}h\to\boxed{\frac25},
\]
using $f(1,1)=0$. (Check: $\grad f(1,1)=\ip{2,-1}$, $\ip{2,-1}\cdot\ip{\frac35,\frac45}=\frac25$.\checkok)

\solhead{8}{7}
$\grad f=\ip{2\cos(2x+3y),\;3\cos(2x+3y)}$; at $(-6,4)$, $2x+3y=0$, so
$\grad f=\ip{2,3}$. $\vu=\ip{\cos(-\frac\pi6),\sin(-\frac\pi6)}=\ip{\frac{\sqrt3}{2},-\frac12}$.
$\Du f=\sqrt3-\frac32=\boxed{\frac{2\sqrt3-3}{2}}\approx0.232$.

\solhead{9}{12}
$\grad f=\ip{2x,4y,6z}$, $\grad f(4,-1,1)=\ip{8,-4,6}$.

\textbf{(a)} Maximum rate $\norm{\ip{8,-4,6}}=\sqrt{116}=\boxed{2\sqrt{29}}$, in the
direction of $\ip{4,-2,3}$.

\textbf{(b)} Check: $16+2+3=21$\checkok. Tangent plane:
$8(x-4)-4(y+1)+6(z-1)=0\iff\boxed{4x-2y+3z=21}$. Normal line:
$x=4+8t$, $y=-1-4t$, $z=1+6t$ (or with direction $\ip{4,-2,3}$).

\solhead{10}{8}
$D_{\vi}f=f_x$ and $D_{\vj}f=f_y$, so $\grad f(P)=\boxed{\ip{-3,4}}$. Fastest decrease:
direction $-\grad f=\ip{3,-4}$, unit $\ip{\frac35,-\frac45}$, at rate $-5$. In the
direction of $\vi+\vj$: $\ip{-3,4}\cdot\frac{1}{\sqrt2}\ip{1,1}=\boxed{\frac{1}{\sqrt2}}=\frac{\sqrt2}{2}$.

\solhead{11}{6}
(i)--R: unit length ($\norm{\ip{-y,x}}=r$), perpendicular to $\vx$, and at $(1,0)$ it
is $\ip{0,1}$: counterclockwise. (ii)--S: $\vx$ itself. (iii)--Q: horizontal
component $0$; vertical component $x$. (iv)--P: constant.

\solhead{12}{8}
$M_y=e^x\cos y=N_x$, continuous on $\R^2$: conservative.
$f=\int(e^x\sin y+2x)\,\dd x=e^x\sin y+x^2+g(y)$;
$f_y=e^x\cos y+g'(y)=e^x\cos y-1\Rightarrow g=-y+K$.
$\boxed{f=e^x\sin y+x^2-y+K}$.

\solhead{13}{8}
Test: $M_y=2yz^3=N_x$; $N_z=6xyz^2=P_y$; $P_x=3y^2z^2=M_z$. Conservative.
$f=\int y^2z^3\,\dd x=xy^2z^3+g(y,z)$; $f_y=2xyz^3+g_y=2xyz^3\Rightarrow g=h(z)$;
$f_z=3xy^2z^2+h'=3xy^2z^2\Rightarrow h=K$. $f(1,1,1)=1+K=4\Rightarrow K=3$:
$\boxed{f=xy^2z^3+3}$.

\solhead{14}{8}
$M_y=1=N_x$ and $N_z=0$ but $P_y=xz$; also $P_x=yz\neq0=M_z$. Not identically equal,
so $\vF$ is \textbf{not conservative}.
\emph{Method breakdown:} $f=xy+g(y,z)$; $f_y=x+g_y=x\Rightarrow g=h(z)$; then
$f_z=h'(z)$ must equal $xyz$, which depends on $x$ and $y$: impossible.
